\clearpage\thispagestyle{empty}
\pdfbookmark[0]{Quellen zu den beiden Briefen an Bode von 1808}{bode1808letters-sources}
\begin{center}\textbf{Quellen zu den beiden Briefen an Bode von 1808}\end{center}
\noindent Diese KI-Übertragungen ins heutige Deutsch folgen den beiden Briefen von Carl Friedrich Gauß an Johann Elert Bode vom 24. Januar und 23. Juli 1808 in \href{https://gdz.sub.uni-goettingen.de/id/PPN236007467}{Gauß’ \textit{Werken}, Bd.~VI} (Göttingen, 1874), S.~293--296. Bode druckte beide Briefe zusammen im \textit{Astronomischen Jahrbuch für das Jahr 1811} (Berlin, 1808), \href{https://www.digitale-sammlungen.de/en/view/bsb10538340?page=147}{S.~135}--\href{https://www.digitale-sammlungen.de/en/view/bsb10538340?page=151}{139}. Im Exemplar der Bayerischen Staatsbibliothek steht der Januarbrief auf \href{https://www.digitale-sammlungen.de/en/view/bsb10538340?page=147}{S.~135}, \href{https://www.digitale-sammlungen.de/en/view/bsb10538340?page=148}{S.~136} und im oberen Teil von \href{https://www.digitale-sammlungen.de/en/view/bsb10538340?page=149}{S.~137}, der Julibrief am Fuß von \href{https://www.digitale-sammlungen.de/en/view/bsb10538340?page=149}{S.~137} sowie auf \href{https://www.digitale-sammlungen.de/en/view/bsb10538340?page=150}{S.~138} und \href{https://www.digitale-sammlungen.de/en/view/bsb10538340?page=151}{S.~139}.
\par\medskip
In der ersten Zeile zum 4. Januar beginnt die Rektaszension in der Werkausgabe mit 325, im Jahrbuch mit 525. Für das Dollond-Instrument bei Thiarks’ Trabantenbeobachtung am 24. Juli gibt die Werkausgabe $3\tfrac14$ f. an, das Jahrbuch $3\tfrac12$ f. Beide Ausgaben nennen den 24. Juli im Brief vom 23. Juli. Die Werkausgabe beschreibt die Ephemeride des Januarbriefs in eckigen Klammern und lässt zwei mit B. gezeichnete Anmerkungen des Jahrbuchs aus. Die Übertragungen folgen der Werkausgabe. KI-Übertragung und KI-Satz.
\clearpage
